\section{Finite-group rigidity at positive error}\label{sec:noisy52}
The exact classification has a robust counterpart at arbitrary width.
The relevant reduction is not to minimum rank, but to the part of the
orthogonal representation on which every nonzero orbit is infinite.
Throughout this section the numerical interface is that of
Section~\ref{sec:lifts51}; in particular, the signed coordinate seeds,
all coordinate queries, and an identity command are present. Write
$G=\langle U_a:a\in\calA\rangle$ and let $H=\overline G\subset O(D)$.
The width remains clocked and nonuniform in the horizon.

\begin{lemma}[The finite-orbit subspace]\label{lem:finite-orbit52}
The set
\[
 F=\{v\in\R^D: Gv\text{ is finite}\}
\]
is a $G$-invariant linear subspace, the image of $G$ on $F$ is finite,
and $E=F^\perp$ is invariant. If $G$ is infinite, then $E\ne0$ and
every nonzero vector of $E$ has an infinite $H$-orbit. Moreover
\begin{equation}\label{eq:seed-moving52}
 s_*:=\max_i\|P_Ee_i\|\ \ge\sqrt{\dim E/D}\ \ge D^{-1/2}.
\end{equation}
\end{lemma}
\begin{proof}
For vectors with finite orbits, the orbit of any linear combination
is contained in the corresponding finite set of linear combinations.
This proves linearity. Invariance follows from $G(gv)=Gv$. Choose a
basis of $F$; the union of its finite orbits is a finite spanning set
permuted by $G$, so the induced image on $F$ is finite. Orthogonality
makes $E$ invariant. If $E=0$, the images of a basis determine only
finitely many matrices, and $G$ is finite. A nonzero vector of $E$
with finite orbit would belong to $F$, a contradiction. Taking the
closure does not change a finite orbit, so the same conclusion holds
for $H$. Finally $\sum_i\|P_Ee_i\|^2=\operatorname{tr}P_E=\dim E$.
\end{proof}

We use normalized Haar measure $\mu$ on the compact group $H$; its
existence and full support are standard \cite{Folland16}. For $E\ne0$
and $k\ge1$ define
\begin{equation}\label{eq:haar-gap52}
 g_E(k)=\min_{u,v_1,\ldots,v_k\in\Sphere(E)}
 \int_H\left(1-\max_j\ip{v_j}{hu}\right)\,d\mu(h),
\end{equation}
where $\Sphere(E)$ is the Euclidean unit sphere of $E$. Restricting
an orthogonal map to $E$ is understood in this and subsequent formulas.

\begin{lemma}[Finite executable distortion on the moving subspace]
\label{lem:haar-word52}
If $G$ is infinite, $g_E(k)>0$ for every finite $k$. There exist an
integer $B\ge1$, a probability law $p$ on actual words of length $B$,
and $d\in(0,1)$ such that
\begin{equation}\label{eq:moving-distortion52}
 \min_{u,v_1,\ldots,v_k\in\Sphere(E)}
 \sum_w p_w\left(1-\max_j\ip{v_j}{U_wu}\right)\ge d.
\end{equation}
One may take $d=g_E(k)/2$. No bound on the time needed to find $B$
is asserted for arbitrary real input.
\end{lemma}
\begin{proof}
The parameter space in \eqref{eq:haar-gap52} is compact and the
integrand is continuous, uniformly in its parameters. The minimum is
therefore attained. If it were zero, nonnegativity and full support
of Haar measure would give $hu\in\{v_1,\ldots,v_k\}$ for every
$h\in H$. This contradicts Lemma~\ref{lem:finite-orbit52}.
The value is at most one: choosing all centers parallel to the
average of $hu$, when nonzero, gives a nonnegative average pairing.
Thus $g_E(k)/2\in(0,1)$.

The closure of the positive word semigroup is $H$. To see this,
compactness of the powers of a fixed orthogonal generator supplies
positive powers tending to the identity; the preceding powers then
tend to its inverse. A finite-order generator has its inverse as an
actual positive power. Thus the semigroup closure contains all
generator inverses and equals the group closure.
For completeness, recurrence of positive powers follows by covering
the compact group by finitely many balls and subtracting the indices
of two powers in the same ball; their quotient tends to the identity.

For $\eta=g_E(k)/2$, choose finitely many positive words whose
matrices are an $\eta$-net of $H$ in operator norm, using a strictly
smaller covering radius if necessary. Partition $H$ into Borel cells,
each assigned to one such matrix within distance $\eta$, and assign
the Haar mass of a cell to its representative word. The functions
$h\mapsto1-\max_j\ip{v_j}{hu}$ are uniformly $1$-Lipschitz.
Replacing Haar measure by this finite law changes every integral by
at most $\eta$. Padding the words with identity letters to a common
length, and increasing that length to at least one, proves
\eqref{eq:moving-distortion52}.
\end{proof}

The endpoint argument from Lemma~\ref{lem:packet} applies on any
invariant Euclidean subspace. We repeat the short calculation here
because its calibration, rather than an assumption about statewise
physical vectors, is decisive.

\begin{theorem}[Noisy occupation at arbitrary width]\label{thm:noisy-budget52}
Suppose $G$ is infinite. Choose a signed coordinate seed $e_i$ with
$\|P_Ee_i\|=s_*$, and put
\begin{equation}\label{eq:kappa-moving52}
 \kappa_E=\frac{\rho s_*}{\sqrt D}-2\varepsilon.
\end{equation}
Assume $\kappa_E>0$. For $k\ge1$, take $B,p,d$ from
Lemma~\ref{lem:haar-word52}, and put $\chi=-\log(1-d)>0$.
Every machine of worst-word total-variation error at most
$\varepsilon$ satisfies
\begin{align}
 \lfloor N/B\rfloor\chi&\le\log(1/\kappa_E)
       &&\text{if its peak is at most }k,\label{eq:noisy-peak52}\\
 \#\{1\le t\le N:K_t\le k\}
       &\le B-1+\frac{B\log(1/\kappa_E)}{\chi}.
       \label{eq:noisy-occupation52}
\end{align}
The unselected registers in the second assertion may have arbitrary
finite width. There is no reachability-rank or conditioning hypothesis.
\end{theorem}
\begin{proof}
Fix the seed $e_i$ and let the target direction at a prefix $u$ be
$Y=U_uP_Ee_i/s_*\in\Sphere(E)$. If the machine state is $S$, set
$z_s=\E[Y\mid S=s]$ and $q=\E\|z_S\|$. Initially $q=1$.
Across a complete independently sampled word with law $p$, composition
of all stochastic updates gives one endpoint kernel $T(j\mid s,w)$.
The fresh random variables and the complete word are independent of
the past. Consequently
\[
 z'_j=\E[U_Wz_S\mid S'=j].
\]
Choose $v_j\in\Sphere(E)$ in the direction of $z'_j$ when this vector
is nonzero. For at most $k$ endpoint labels, padding the centers when
necessary, one obtains
\begin{align*}
 q'&=\sum_s\Pp(S=s)\sum_w p_w\sum_jT(j\mid s,w)
                                      \ip{v_j}{U_wz_s}\\
   &\le\sum_s\Pp(S=s)\|z_s\|
          \max_{u\in\Sphere(E)}\sum_wp_w\max_j\ip{v_j}{U_wu}
    \le(1-d)q.
\end{align*}
States with $z_s=0$ contribute zero. This reasoning uses every actual
hidden state, not a projected or rank-tight substitute. A deterministic
word cannot increase $q$, by orthogonality and conditional Jensen.

Let $D(S_N)\in[-1,1]^D$ be the vector of terminal decoder means.
For each complete word the realized mean vector differs from
$\rho U_we_i$ by at most $2\varepsilon$ per coordinate. Since
$Y_N=U_wP_Ee_i/s_*$ and $F\perp E$ are invariant,
\[
 \E\ip{Y_N}{D(S_N)}\ge\rho s_*-2\sqrt D\,\varepsilon.
\]
On the other hand, conditioning on $S_N$ bounds the left side by
$\sqrt D\,q_N$. Hence $q_N\ge\kappa_E$. If $J$ disjoint intervals
carry independent copies of $p$, their endpoint contractions give
$\kappa_E\le(1-d)^J$. The machine must pass this law because its
correctness is wordwise; the law is only a test, not a change in the
input model. Consecutive intervals prove \eqref{eq:noisy-peak52}.

For the occupation bound discard endpoints before $B$, choose the
first remaining endpoint, and discard subsequent candidates at
strictly smaller distance than $B$; repeat. If $M$ is the original
number of endpoints, this gives $M-(B-1)\le BJ$ disjoint intervals.
Their start registers and their intermediate registers are unrestricted.
Combining with the calibration gives \eqref{eq:noisy-occupation52}.
\end{proof}

\begin{theorem}[Positive-error finite-group classification]
\label{thm:noisy-finitegroup52}
For every finite orthogonal alphabet with an identity, signed
coordinate seeds, all coordinate queries, and $0<\rho<1$, fix
\begin{equation}\label{eq:noise-range52}
                    0\le\varepsilon<\rho/(2D).
\end{equation}
Then the following are equivalent: the clocked peaks $W_{N,\varepsilon}$
are bounded over all horizons; the generated group is finite; and
one finite stationary permutation machine realizes the experiment
exactly at every horizon. If the group is infinite, then
$W_{N,\varepsilon}\to\infty$, with the arbitrary-width occupation
bound of Theorem~\ref{thm:noisy-budget52}.
For a particular infinite action, \eqref{eq:noise-range52} may be
replaced by $\varepsilon<\rho s_*/(2\sqrt D)$.
\end{theorem}
\begin{proof}
For an infinite group, \eqref{eq:seed-moving52} and
\eqref{eq:noise-range52} imply $\kappa_E>0$. A purported uniform
peak $k$ contradicts \eqref{eq:noisy-peak52} as $N\to\infty$.
For a finite group, the orbit permutation construction in
Theorem~\ref{thm:finitegroup51} has zero error and a fixed finite
number of labels. Finally, $W_{N,\varepsilon}$ is nondecreasing in
$N$: replace a terminal decoder $d_j$ of an $(N+1)$-horizon machine
by $T_{N+1,I}d_j$. This gives a legal $N$-horizon decoder with the
same error bound. Unboundedness therefore implies divergence.
\end{proof}

The error interval in this theorem is sufficient, not claimed sharp.
For $\varepsilon\ge\rho/2$, the fair binary answer already realizes
every target within tolerance, irrespective of $G$. The intermediate
noise regime is not classified here. A fixed subspace must not be
ignored: for $\operatorname{diag}(R_\theta,1)$ the distortion computed
on all of $\R^3$ vanishes, whereas the finite-orbit reduction removes
the fixed axis and leaves the nontrivial planar occupation bound.
Thus Theorem~\ref{thm:noisy-finitegroup52} is not obtained by simply
asserting that the old full-space distortion is positive.

\subsection{An explicit arithmetic rate without a Diophantine hypothesis}
The finite word proof admits a direct quantitative certificate in a
broad elementary class. Its rate is deliberately coarse; its input
is rational height, not an uncomputed variational minimum.

\begin{theorem}[Rational-height occupation]\label{thm:height52}
In dimension two suppose the alphabet contains $I$ and a rotation
with eigenvalue $\lambda=(a+ib)/c$, where $a,b,c\in\Z$, $c\ge2$,
$a^2+b^2=c^2$, and $\lambda$ is not a root of unity. Other commands
may be present and need not commute. For
$\kappa=\rho/\sqrt2-2\varepsilon>0$, any feasible peak $K\ge1$
satisfies
\begin{equation}\label{eq:height-peak52}
 N<2K\{1+16c^{4K}\log(1/\kappa)\}.
\end{equation}
More generally,
\begin{equation}\label{eq:height-occupation52}
 \#\{1\le t\le N:K_t\le K\}
       \le2K-1+32Kc^{4K}\log(1/\kappa).
\end{equation}
In particular,
\[
 W_{N,\varepsilon}\ge
       \frac{\log N-O(\log\log N)}{4\log c}
       \quad(N\to\infty),
\]
where $\rho,\varepsilon,c$ are fixed. This applies to the rational
noncommuting family of Corollary~\ref{cor:rational51} with $c=5$.
\end{theorem}
\begin{proof}
Use the uniform law on the $2K$ words with products
$I,R,\ldots,R^{2K-1}$, each padded to length $B=2K$.
For $1\le q\le2K-1$, the nonzero Gaussian integer
$(a+ib)^q-c^q$ has modulus at least one. Hence
$|\lambda^q-1|\ge c^{-q}\ge c^{-2K}$.
For every starting unit vector, distinct points of this word orbit
are separated by at least $c^{-2K}$. An open Euclidean cap of radius
$\eta=c^{-2K}/2$ therefore contains at most one point and has mass
at most $1/(2K)$. The union of $K$ such caps has mass at most $1/2$.
Outside it, the distortion is at least $\eta^2/2$, giving the
uniform finite-word gap
\[
                         d\ge c^{-4K}/16.
\]
This is exactly the cap estimate of Proposition~\ref{prop:caps},
now with a height-certified separation. Apply the endpoint proof
above with $E=\R^2$ and $\chi\ge d$. It gives
$\lfloor N/(2K)\rfloor\le16c^{4K}\log(1/\kappa)$ and the stated
occupation bound. Taking logarithms proves the asymptotic assertion:
when $K$ exceeds a constant times $\log N$ it is immediate, and
otherwise $\log K=O(\log\log N)$ in \eqref{eq:height-peak52}.
\end{proof}

The bound controls hidden labels, not the vertices of a projected
section, so there is no intervening exponential loss through $V_K$.
The sharper arithmetic laws in the retained appendices concern more
structured word orbits and are not replaced by this height bound.
